\section{GNS Construction}\label{sctn:gns-construction}
Before examining Axiom 2 (Isotony), we need to take a moment to review how AQFT is operationalized.

Operationally, the (abstract) C*-algebras $\mathfrak{U}(\mathbf{B})$ aren't used directly. Instead they are operationalized through the ``GNS construction'', which we now describe.

An \textit{AQFT state}, sometimes simply called a \textit{state}, is an element $\omega$ of the dual space $\mathfrak{U}(\mathbf{B})^*$ that is

\begin{itemize}
    \item \textit{Positive} - for any $a \in \mathfrak{U}(\mathbf{B})$ we have $0 \le \omega(a^*a)$ and
    \item \textit{Normalized} - the operator norm satisfies $\| \omega \|=1$.
\end{itemize}

Furthermore, an AQFT state $\omega$ is said to be \textit{faithful} if for any non-zero $a$ in $\mathfrak{U}(\mathbf{B})$ it follows that $0 < \omega(a^*a)$.

The \textit{GNS construction} is a procedure that given an AQFT state $\omega$ on an (abstract) C*-algebra $\mathfrak{U}(\mathbf{B})$ with a unit, produces a triple $(\mathcal{H}_\omega, \pi_\omega, \Omega_\omega)$ consisting of a Hilbert space $\mathcal{H}_\omega$, a *-homomorphism (and thus representation) $\pi_\omega$ of $\mathfrak{U}(\mathbf{B})$ on the Hilbert space, and $\Omega_\omega$ a distinguished vector in the Hilbert space. (Note\footnote{Technically, one only requires the existence of a ``bounded, approximate identity'' and not a unit to produce a distinguished vector $\Omega_\omega$. However, physical motivation for the existence of a ``bounded, approximate identity'' as opposed to a unit is lacking.}, without the unit in $\mathfrak{U}(\mathbf{B})$, the GNS construction is unable to produce the distinguished vector $\Omega_\omega$.)

Furthermore, the GNS construction's representation $\pi_\omega$ is faithful if the AQFT state $\omega$ one starts with is faithful. (This justifies the name ``faithful'' being applied to an AQFT state.)

It is in the Hilbert space $\mathcal{H}_\omega$ that ``normal physics'' occurs. In other words elements of $\mathcal{H}_\omega$ are what one thinks of as normal QFT states, and self-adjoint elements of $\pi_\omega(\mathfrak{U}(\mathbf{B}))''$ (the von Neumann algebra generated by $\pi_\omega(\mathfrak{U}(\mathbf{B}))$) are interpreted as the normal observables of QFT\footnote{Given $\mathcal{N}$ a self-adjoint subset of $\mathcal{B}(\mathcal{H}_\omega)$ the set of bounded operators on $\mathcal{H}_\omega$ ( i.e. $\mathcal{N} \subseteq \mathcal{B}(\mathcal{H}_\omega)$ such that $\mathcal{N}^*=\mathcal{N}$ ) we define $\mathcal{N}' \equiv \{ N' \in \mathcal{B}(\mathcal{H}_\omega) : NN' - N'N = 0 \; \forall N \in \mathcal{N} \}$ and similarly define the shorthand $\mathcal{N}'' \equiv (\mathcal{N}')'$.}. Finally, the distinguished vector $\Omega_\omega$ is interpreted as being a vacuum state.

Also note that this GNS construction can in addition be applied to $\mathfrak{U}$ the algebra of quasilocal observables that appears in \textbf{Axiom 4} \textit{(Quasilocal Algebra)}. In other words one can have an AQFT state $\omega$ that is an element of $\mathfrak{U}^*$ and use this AQFT state and the GNS construction to create the triple $(\mathcal{H}_\omega, \pi_\omega, \Omega_\omega)$. (Again, without a unit in $\mathfrak{U}$, the GNS construction is unable to produce the distinguished vector $\Omega_\omega$.)
